\chapter{GİRİŞ}
\label{ch:giris}

% --- 1.1 ---
\section{Otonom Ajanlar ve Modern Yapay Zeka Mimarileri}
\label{sec:ajan-mimari}

Modern yapay zeka uygulamalarının temelinde, devasa metin ve kod veri kümeleri üzerinde eğitilmiş Büyük Dil Modelleri (LLM) yer almaktadır. Bu modeller, kendilerine sunulan bilgiyi anlama ve üretme konusunda yüksek kapasiteye sahip olsalar da, tek başlarına karmaşık ve çok aşamalı görevleri yerine getirmede sınırlı kalmaktadır. Bu noktada ajan tabanlı sistemler devreye girmektedir. Bir ajan, LLM'yi karar alma çekirdeği olarak kullanır; hedefe yönelik planlama yapar, dosya okuma veya kod yürütme gibi yardımcı araçları etkinleştirir ve çevreden aldığı geri bildirime göre stratejisini günceller.

Örneğin finans teknolojileri alanında çalışan bir yazılım mühendisinin, on binlerce dosyadan oluşan yeni bir ödeme altyapısına katıldığını düşünelim. Geliştiricinin mimariyi anlaması, iş akışlarını kavraması ve basit bir değişikliği güvenle uygulayabilmesi çoğu zaman haftalar almaktadır. İnsan gözetiminde yürüyen bu uzun uyum sağlama süreci, kod tabanından yapılandırılmış bilgi çıkarabilen otonom sistemlerin yokluğunda daha da uzamakta; iş devri ve bakım maliyetleri yükselmektedir.

Bu mimarilerin en belirgin teknik kısıtlaması, LLM'lerin "bağlam penceresi" olarak bilinen veri işleme sınırıdır. Bir model aynı anda yalnızca sınırlı miktarda bilgiyi hafızasında tutabilir. Dolayısıyla kapsamlı bir yazılım projesinin tamamını tek seferde değerlendirmek geleneksel ajan yaklaşımlarıyla mümkün olamamaktadır.

% --- 1.2 ---
\section{Yazılım Mühendisliğinde Ajan Uygulamaları}
\label{sec:ajan-uygulamalari}

Ajan tabanlı sistemler, yazılım geliştirme yaşam döngüsünün (SDLC) birçok aşamasını otomatize etme potansiyeline sahiptir. Bu sistemler, basit kod tamamlama araçlarının çok ötesine geçerek aktif roller üstlenmeye başlamıştır. Güncel uygulama alanları arasında şunlar bulunmaktadır:

\begin{itemize}
    \item \textbf{Otomatik Hata Ayıklama (Debugging):} Bir kod tabanındaki hataları tespit etmek için test senaryoları yürütebilen, günlük (log) dosyalarını analiz edebilen ve hatanın kaynağını belirleyebilen ajanlar.
    \item \textbf{Talep Üzerine Kod Üretimi:} "Kullanıcılar için bir giriş (login) sayfası oluştur" gibi doğal dil taleplerini alıp, gerekli dosya yapılarını ve kod bloklarını üreten sistemler.
    \item \textbf{Mevcut Kodun Yeniden Yapılandırılması (Refactoring):} Verimsiz veya eski kod standartlarına göre yazılmış bölümleri analiz edip, daha performanslı veya modern eşdeğerleriyle güncelleyen ajanlar.
    \item \textbf{Kod Analizi ve Dokümantasyon:} Bu projenin de odak noktası olan, mevcut bir kod tabanını inceleyerek sistemin mimarisini, veri akışlarını ve bileşenlerin sorumluluklarını anlatan teknik dokümantasyon üreten sistemler.
\end{itemize}

% --- 1.3 ---
\section{Projenin Amacı}
\label{sec:proje-amaci}

Yazılım geliştirme süreçlerinde, "geniş ve yabancı bir kod tabanını anlamak, yazılım geliştiricilerin karşılaştığı en önemli zorluklardan biridir". Yeni bir geliştiricinin bir projeye dahil olması  veya tecrübeli bir geliştiricinin projenin daha önce dokunmadığı bir bölümüne müdahale etmesi gerektiğinde, mevcut dokümantasyon genellikle eksik, güncel olmayan veya yüzeysel kalmaktadır. Bu durum, projelere adaptasyon sürecinin haftalar, hatta aylar sürmesine neden olarak ciddi bir zaman ve verimlilik kaybına yol açmaktadır.

Tıpkı geniş ölçekli çevrim içi platformlarda olduğu gibi, modern bir kurumsal kod tabanı binlerce dosya, on binlerce işlev ve milyonlarca satır kod içerebilir. Bu büyüklükteki ve yüksek derecede bağlı verinin, tek bir kişinin kısa sürede incelemesi ve anlamlandırması mümkün değildir. Son yıllarda bu sorunu aşmak için büyük dil modellerinden yararlanan ajan tabanlı yapay zekâ çözümleri geliştirilmektedir. Ancak söz konusu çözümler bile yüksek hacimli kod yığınını bütünüyle kavramada bağlam penceresi sınırlamasına takılmaktadır.

Bu sınırlamanın temel nedeni, LLM'lerin belirli bir bağlam uzunluğunu aşamaması ve bu nedenle kod tabanındaki anlamsal bütünlüğü tek seferde kuramamasıdır. Kod tabanını bütünüyle modele yüklemek ya da aynı ajanın tüm dosyaları art arda okumasını beklemek, bağlamın dağılmasına ve yapılan değerlendirmelerin güvenilirliğinin azalmasına yol açmaktadır.

Bu çalışma kapsamında, bağlam penceresi kısıtlamasını aşmak için \textbf{alt-ajan mimarisi} adını verdiğimiz özgün bir yöntem benimsenmiştir. Söz konusu yaklaşımda kapsamlı analiz görevi daha küçük ve yönetilebilir alt görevlere bölünmektedir.

Projenin temel hedefi, söz konusu mimariyi kullanarak bir "yapay zekâ eğitmeni" prototipi geliştirmektir. Sistem iki ana aşamada işlemektedir:
\begin{enumerate}
    \item \textbf{Aşama 1: Bilgi Bankası Üretimi:} Ana ajan, analiz edilecek yazılım projesini mantıksal modüllerine ayırır. Her bir bölümün analizi için, o bölüme özel, izole bir çalışma alanına ve kısıtlı araçlara sahip bir alt-ajan görevlendirir. Alt-ajan kendi analizini tamamlar ve bulgularını (markdown dosyaları, diyagramlar) belirlenen dizin altında merkezi bir bilgi bankası oluşturacak şekilde kaydeder. Bu yaklaşım, ana ajanın bağlamını temiz tutar ve görevi ölçeklenebilir hale getirir.
    \item \textbf{Aşama 2: Eğitim Dokümanı Üretimi:} İlk aşamada üretilen ve kod tabanının bütününü özetleyen bilgi bankası bu bölümde girdi olarak kullanılır. Eğitim üretici bileşeni, söz konusu bilgi birikimini ve doğrudan kaynak koddan alınan güncel örnekleri bir araya getirerek geliştiricilere yönelik öğretici dokümanlar üretir.
\end{enumerate}

    Bu projenin seçilmesinin temel gerekçesi, özellikle üniversite öğrencileri ve stajyer geliştiriciler için kaynak hazırlanırken karşılaşılan bilgi derleme güçlüğünü ortadan kaldırmaktır. Klasik dokümantasyon süreçleri uzun zaman aldığından pek çok kurum güncel eğitim materyali sunamamaktadır; alt-ajan destekli yarı otonom bir çözüm bu boşluğu kapatmayı amaçlamaktadır.

Çalışmanın özgün katkıları aşağıdaki başlıklar altında toplanabilir:
\begin{itemize}
        \item \textbf{Kararlı görev bölme yaklaşımı:} Ana ajanın kod tabanının dizin yapısını inceleyerek ilk seviye klasörlere göre tutarlı alt görevler oluşturması sağlanmış, böylece karmaşık planlama adımlarının yalın bir kural ile yürütülmesi mümkün kılınmıştır.
        \item \textbf{İzole çalışma alanları:} Her alt-ajanın yalnızca kendi çalışma dizinine yazabildiği, kod tabanına ise salt okunur eriştiği güvenli bir görev yürütme modeli tasarlanmış, bağlam taşması ve veri bütünlüğü riskleri azaltılmıştır.
        \item \textbf{İki aşamalı üretim hattı:} Bilgi bankası ile başlayan ve eğitim dokümanı üretimiyle tamamlanan süreç, dokümantasyon kalitesini yükseltirken ortaya çıkan içeriğin doğrudan öğretim amaçlı kullanılmasına olanak tanımaktadır.
\end{itemize}

% -------------------------------------------------------------------
% BÖLÜM 2
% -------------------------------------------------------------------

\chapter{ÖN İNCELEME}
\label{ch:on-inceleme}


% --- 2.1 ---
\section{Projeye Olan İhtiyaç}
\label{sec:ihtiyac}

Yazılım endüstrisinde kapsamlı ve yabancı kod tabanlarının anlaşılması, özellikle bilgi birikiminin kurumsal hafızada yeterince belgelenmemesi nedeniyle önemli bir darboğaz oluşturmaktadır. Mevcut uygulamada yeni bir geliştirici, kodun işlevsel bölümlerini inceleyebilmek için çok sayıda dosyayı tek tek okumakta; bu süreç hem yüksek zaman maliyeti doğurmakta hem de yorumlayanın bilgi düzeyine bağlı olarak tutarsız sonuçlar üretmektedir. 

Bu problem alanında öne çıkan ilk yaklaşım, tek bir ajanla çalışan kod inceleme çözümleridir. AutoGPT ve benzeri erken dönem örnekler, tek bir büyük dil modelinin (LLM) araçlarla desteklenmesi sayesinde belirli görevleri yerine getirebilmiştir. Ancak bu tür sistemler, bağlam penceresi sınırlaması nedeniyle çok sayıda dosyadan aynı oturum içinde sağlıklı çıkarım yapamamış, ayrıca uzun görevleri yürütürken odak kaybı yaşamıştır. 

Sonraki aşamada geliştirilen yönetici-işçi düzenindeki çoklu ajan çözümleri ise bir yöneticinin ardışık görevler atadığı yapıları benimsemiştir. Langchain gibi çatılar kullanılarak geliştirilen yönetici-işçi örnekleri bu yaklaşıma bir örnektir. Bu sistemler iş yükünü paylaştırsa da, işçiler tarafından üretilecek içeriklerin yönetici aracılığıyla toplanması yüksek hacimli analizlerde yine bağlam taşmasına yol açmaktadır. Ayrıca her işçinin yalnızca yöneticinin verdiği sırayla çalışabilmesi, dinamik ve geniş kod tabanlarında esnekliği kısıtlamaktadır. 

Bu inceleme, mevcut yaklaşımların zayıf yönlerini ortaya koyarak alt-ajan mimarisinin gerekliliğini vurgulamaktadır. Tek ajanlı sistemlerin bağlam penceresi ve uzun soluklu planlama sorunları ile yönetici-işçi mimarilerinin verimlilik ve esneklik kısıtları, bu çalışmada önerilen alt-ajan yaklaşımının temel motivasyonunu oluşturmaktadır.

% --- 2.2 ---
\section{Proje Kapsamı}
\label{sec:kapsam}

Literatürdeki yaklaşımlar ve mevcut sınırlılıklar doğrultusunda, projenin kapsamı aşağıdaki başlıklar altında tanımlanmıştır:

\begin{itemize}
    \item \textbf{Girdi kod tabanının incelenmesi:} Sistem, \texttt{smolagents} hızlı başlangıç şablonu gibi bir temel üzerinde, sağlanan depo ve dizin yapısını güvenli dosya erişimi yöntemleriyle tarayacaktır. Projenin temel bileşenleri ve modülleri belirlenerek ana ajanın görev planı oluşturulacaktır.
    \item \textbf{Alt-ajan destekli bilgi bankası üretimi:} Her alt-ajanın yalnızca yetkilendirildiği dizinleri okuyabildiği ve kendi çalışma klasörüne yazı kaydedebildiği bir güvenlik düzeni kurulacaktır. Alt-ajanlar, sorumlu oldukları modüllerin mimarisini, veri akışını ve önemli fonksiyonlarını ayrıntılandıran teknik belgeler oluşturacaktır.
    \item \textbf{Eğitim odaklı çıktı üretimi:} Bilgi bankasında toplanan içerik, projenin ana hedefi doğrultusunda, bir eğitim materyaline dönüştürülecektir. Alt-ajanlardan gelen raporlar, geliştiricilerin sisteme hızlı uyum sağlamasını destekleyecek bir anlatım diline dönüştürülecek; kod parçaları ve diyagramlarla desteklenecektir.
    \item \textbf{Çevrim içi ve çevrim dışı kullanım senaryoları:} Üretilen materyallerin yalnızca yeni başlayanlar için değil, denetim veya bakım görevleri için de başvuru kaynağı olabilmesi hedeflenmektedir. Bu nedenle çıktılar dizin bazında düzenlenmiş, yeniden kullanılabilir belgelerden oluşacaktır.
\end{itemize}

% --- 2.3 ---
\section{Projenin Gereksinimleri}
\label{sec:gereksinimler}

Sistemin hedeflenen işlevleri yerine getirebilmesi için öngörülen gereksinimler aşağıda sunulmuştur:

\begin{itemize}
    \item \textbf{Girdi kod tabanı:} Analiz edilecek depo, orta ölçekli bir Python projesi olmalı; dizin yapısı ana ajan tarafından okunabilir biçimde bulunmalıdır.
    \item \textbf{Model altyapısı:} Alt-ajanlar ve ana ajan, \texttt{LiteLLM} aracılığıyla erişilen ve araç çağırma kabiliyeti destekleyen bir büyük dil modeli üzerinde çalışacaktır. Tercihen Google Gemini ailesi veya eşdeğer bir model kullanılacaktır.
    \item \textbf{Geliştirme ortamı:} Python 3.10 ve üzeri sürüm, \texttt{smolagents} çatısı, güvenli dosya erişimini sağlayan yardımcı kütüphaneler ve otomatik testleri çalıştıracak altyapı gereklidir.
    \item \textbf{Yardımcı araçlar:} Kod tabanından güvenli okuma, alt-ajan çalışma alanlarına yazma, dizin ağacını listeleme ve alt-ajan çıktılarının birleştirilmesi için özelleştirilmiş araç fonksiyonları sağlanacaktır.
    \item \textbf{Pedagojik çıktılar:} Bilgi bankası ve eğitim materyali üretimi için markdown düzenleyicileri, diyagram oluşturma araçları ve gerektiğinde görsel temsil sağlayacak bileşenler kullanılacaktır.
\end{itemize}